\documentclass{article}
\usepackage{iclr2027_conference,times}
\usepackage[T1]{fontenc}
\usepackage{lmodern}
\input{math_commands.tex}
\usepackage{amsmath,amssymb,amsthm,mathtools,bm,booktabs,microtype,algorithm,algorithmic}
\usepackage{hyperref}
\usepackage{url}

\newtheorem{theorem}{Theorem}
\newtheorem{lemma}[theorem]{Lemma}
\newtheorem{corollary}[theorem]{Corollary}
\newtheorem{proposition}[theorem]{Proposition}
\newtheorem{definition}[theorem]{Definition}
\newtheorem{remark}[theorem]{Remark}
\newcommand{\W}{\mathcal W}
\newcommand{\Reg}{\operatorname{Reg}}
\newcommand{\DReg}{\operatorname{DReg}}
\newcommand{\ind}{\mathbf 1}

\title{No Price for an Unknown Discount at a Known Deadline}

\author{\textbf{Pahan Dewasurendra} \\ Johns Hopkins University}

\hypersetup{pdftitle={No Price for an Unknown Discount at a Known Deadline},pdfauthor={Pahan Dewasurendra},pdfkeywords={price, unknown, discount, known, deadline, discounted, regret, evaluates, online},pdfsubject={preprint}}
\begin{document}
\maketitle

\begin{abstract}
Discounted regret evaluates an online learner with exponentially larger weight on recent losses. If the discount factor is known, online gradient descent attains regret of order $DG/\sqrt{1-\lambda}$. A recent result adapts to an unknown discount factor over a continuum at an extra factor $\sqrt{\log T}$, using a logarithmic collection of experts and discounted meta-algorithms. We show that this adaptivity cost vanishes when the terminal horizon $T$ is known. One projected-gradient iterate with the increasing step sizes $\eta_t=D/(G\sqrt{2(T-t+1)})$, chosen without $\lambda$, satisfies simultaneously for every $\lambda\in[0,1]$
\[
\DReg_T(\lambda)\le \sqrt2DG\sqrt{\frac{1-\lambda^T}{1-\lambda}},
\]
where the ratio is $T$ at $\lambda=1$. For every fixed $\lambda$, even an algorithm told $\lambda$ has minimax expected regret at least $1/(4\sqrt2)$ of this upper bound. Thus a single discount-oblivious trajectory is pointwise minimax for the entire continuum, including effective horizons from one to $T$. The mechanism is a terminal-suffix principle. Increasing step sizes make every suffix ending at the known deadline regret-optimal. Abel summation expresses any exponential discount as a probability mixture of these suffixes, and Jensen's inequality produces the exact finite-horizon effective scale. The same argument applies to mirror descent and to every nonincreasing terminal weighting kernel. Besides closing the logarithmic gap, the result reduces the number of learner states and per-round updates from $O(\log T)$ to one.
\end{abstract}

\section{Introduction}

Online learning often discounts stale observations. At a terminal time $T$, the $\lambda$-discounted regret of decisions $w_1,\ldots,w_T$ is
\begin{equation}
\DReg_T(\lambda)=\sum_{t=1}^T\lambda^{T-t}f_t(w_t)-\min_{u\in\W}\sum_{t=1}^T\lambda^{T-t}f_t(u),
\label{eq:dreg}
\end{equation}
where $\lambda\in[0,1]$. The effective number of rounds is
\begin{equation}
H_T(\lambda):=\sum_{j=0}^{T-1}\lambda^j=
\begin{cases}
(1-\lambda^T)/(1-\lambda),&\lambda<1,\\
T,&\lambda=1.
\end{cases}
\label{eq:effective}
\end{equation}
Small $\lambda$ asks for good performance on the recent past. The endpoint $\lambda=1$ recovers ordinary static regret.

Discounted regret has a long history in prediction with expert advice and nonstationary learning~\citep{freund2008hedging,kapralov2011prediction,chernov2010prediction}. It has recently reappeared in online convex optimization (OCO), online conformal prediction, online regression, dynamic regret, and analyses of adaptive optimizers~\citep{zhang2024discounted,jacobsen2024online,ahn2024understanding,xie2026dynamic}. For a known constant discount, projected online gradient descent (OGD) with a tuned constant step size achieves $O(DG\sqrt{H_T(\lambda)})$ regret on a diameter-$D$ domain with $G$-bounded subgradients. The dependence on $H_T(\lambda)$ is minimax optimal~\citep{zhang2024discounted}.

The appropriate discount is rarely known when decisions are generated. Yang et al.~\citep{yang2026uniform} recently gave the first OCO guarantee uniform over a continuous interval of unknown factors. Their Smoothed OGD construction discretizes the factors, runs $O(\!\log T)$ OGD experts, and combines them sequentially with discounted normal predictors. It guarantees
\begin{equation}
\DReg_T(\lambda)=O\!\left(DG\sqrt{\frac{\log T}{1-\lambda}}\right)
\label{eq:prior}
\end{equation}
on its target interval. The extra $\sqrt{\log T}$ is described there as the cost of adapting to the discount. The construction and analysis are substantial because experts tuned to different discounts are evaluated by different weighted objectives.

We prove that no adaptivity price is necessary at a known deadline. The algorithm is one ordinary projected-gradient trajectory. It knows $T$, $D$, and $G$, but it is not given $\lambda$ and does not maintain candidate factors. Its step size increases as the deadline approaches:
\begin{equation}
\eta_t=\frac{D}{G\sqrt{2(T-t+1)}}.
\label{eq:schedule}
\end{equation}
The resulting decisions obey the following deterministic guarantee for every loss sequence:
\begin{equation}
\forall\lambda\in[0,1],\qquad
\DReg_T(\lambda)\le\sqrt2DG\sqrt{H_T(\lambda)}.
\label{eq:headline}
\end{equation}
The universal quantifier is literal. The discount may be selected after observing the entire trajectory. The same decisions attain the optimal order for all effective horizons between one and $T$.

The proof reverses a familiar learning-rate idea. A decreasing schedule gives useful prefix guarantees because the horizon elapsed so far grows. At a fixed deadline, the relevant recency window shrinks as time advances. The increasing schedule in~\eqref{eq:schedule} makes OGD optimal on every \emph{suffix} $[s,T]$ at once:
\begin{equation}
\sum_{t=s}^T(f_t(w_t)-f_t(u))\le\sqrt2DG\sqrt{T-s+1}.
\label{eq:suffixintro}
\end{equation}
An Abel identity then writes exponentially weighted regret as a convex combination of these suffix regrets. The combination puts mass $(1-\lambda)\lambda^{j-1}$ on a suffix of length $j$ and mass $\lambda^{T-1}$ on the full horizon. Its mean suffix length is exactly $H_T(\lambda)$. Jensen's inequality turns~\eqref{eq:suffixintro} into~\eqref{eq:headline} without changing the constant.

This upper bound is pointwise minimax. On a one-dimensional domain, independent Rademacher linear losses give, for every fixed $\lambda$ and every online algorithm,
\begin{equation}
\E[\DReg_T(\lambda)]\ge\frac14DG\sqrt{H_T(\lambda)}.
\label{eq:lowerintro}
\end{equation}
The lower bound grants the learner both $T$ and $\lambda$. Comparing~\eqref{eq:headline} and~\eqref{eq:lowerintro}, the discount-oblivious trajectory is within a factor $4\sqrt2$ of the oracle minimax value at every factor.

The suffix-mixture argument is not specific to Euclidean projection or exponential discounting. We prove the same suffix profile for online mirror descent whenever the Bregman diameter is bounded. We also show that any nonincreasing terminal kernel is a nonnegative mixture of suffix indicators. If its newest loss has unit weight and its weights sum to $H$, the same trajectory has weighted regret at most $\sqrt{2}DG\sqrt H$. The geometric case is special because a matching lower bound has the same effective scale.

\paragraph{Contributions.} Our main results are: (i) one discount-oblivious OGD trajectory with the exact uniform finite-horizon bound~\eqref{eq:headline}, (ii) a pointwise minimax lower bound for every discount including both endpoints, (iii) a terminal-suffix mixture principle for arbitrary nonincreasing kernels, (iv) an online mirror-descent extension, and (v) a reduction from $O(\log T)$ expert states and meta-updates to one state and one projection per round. The deadline is essential to our construction. We separate this fixed-terminal objective from guarantees required at every unknown stopping time.

\section{Setting and comparison}
\label{sec:setting}

Let $\W$ be a nonempty closed convex subset of a Euclidean space with diameter at most $D$. In round $t$, the learner selects $w_t\in\W$. The environment then reveals a convex loss $f_t:\W\to\mathbb R$ and a subgradient $g_t\in\partial f_t(w_t)$ satisfying $\|g_t\|_2\le G$. The loss sequence may be adaptive to past decisions. The learner knows a terminal horizon $T$ before round one. We assume $D,G>0$, since either zero case is immediate.

For a comparator $u\in\W$, define the discounted linearized regret
\[
R_T(\lambda,u):=\sum_{t=1}^T\lambda^{T-t}\bigl(f_t(w_t)-f_t(u)\bigr).
\]
Convexity gives the standard reduction to subgradients, but our final guarantee is stated directly for the convex losses. Since $\DReg_T(\lambda)=\max_{u\in\W}R_T(\lambda,u)$, a bound uniform in $u$ controls~\eqref{eq:dreg}.

Table~\ref{tab:comparison} compares the closest worst-case guarantees. The known-factor row follows from tuned constant-step OGD~\citep{zhang2024discounted,yang2026uniform}. The prior uniform row is Smoothed OGD~\citep{yang2026uniform}. Constants and the lower endpoint imposed there are omitted to emphasize the rates. Our bound holds on the larger closed interval $[0,1]$ and retains the exact finite-horizon effective length $H_T(\lambda)$.

\begin{table}[t]
\centering
\small
\begin{tabular}{@{}lll@{}}
\toprule
Method & Uses $\lambda$ & Regret at $\lambda$ \\
\midrule
Tuned constant-step OGD & yes & $O(DG\sqrt{H_T(\lambda)})$ \\
Smoothed OGD~\citep{yang2026uniform} & no & $O(DG\sqrt{H_T(\lambda)\log T})$ \\
Deadline OGD, this work & no & $\sqrt2DG\sqrt{H_T(\lambda)}$ \\
Oracle lower bound, this work & yes & $\frac14DG\sqrt{H_T(\lambda)}$ \\
\bottomrule
\end{tabular}
\caption{Uniform adaptation at a known deadline. Deadline OGD keeps one vector and makes one projection per round, while Smoothed OGD uses $O(\!\log T)$ expert and combiner states. The lower bound applies separately at each $\lambda$ and allows the algorithm to know that factor.}
\label{tab:comparison}
\end{table}

The objective differs from adaptive or strongly adaptive regret, which asks for control on all intervals and usually uses restarting or aggregation~\citep{daniely2015strongly,cutkosky2020parameter}. We need all intervals with a common right endpoint $T$. This one-sided structure is exactly what lets a single increasing schedule work. Discounted regret with factors revealed online is another distinct problem. Zhang et al.~\citep{zhang2024discounted} give adaptive FTRL guarantees when the factor sequence is supplied to the learner. Here the constant factor is never revealed because every factor is certified simultaneously.

\section{One trajectory for every discount}
\label{sec:upper}

Algorithm~\ref{alg:deadline} is projected OGD~\citep{zinkevich2003online,hazan2016introduction} with the reverse-horizon schedule~\eqref{eq:schedule}. The word ``reverse'' refers only to the step-size clock. Decisions remain causal and updates use no future loss information.

\begin{algorithm}[h]
\caption{Deadline OGD}
\label{alg:deadline}
\begin{algorithmic}[1]
\REQUIRE Convex domain $\W$, diameter bound $D$, gradient bound $G$, deadline $T$
\STATE Choose any $w_1\in\W$
\FOR{$t=1,\ldots,T$}
\STATE Predict $w_t$, observe $g_t\in\partial f_t(w_t)$
\STATE Set $\eta_t=D/(G\sqrt{2(T-t+1)})$
\STATE Set $w_{t+1}=\Pi_{\W}(w_t-\eta_tg_t)$
\ENDFOR
\end{algorithmic}
\end{algorithm}

We begin with the property that does all of the work.

\begin{lemma}[Every terminal suffix]
\label{lem:suffix}
For every comparator $u\in\W$ and every starting time $s\in[T]$, Algorithm~\ref{alg:deadline} satisfies
\[
\sum_{t=s}^T\bigl(f_t(w_t)-f_t(u)\bigr)\le\sqrt2DG\sqrt{T-s+1}.
\]
The inequalities hold simultaneously for all $s$ and $u$ on every loss sequence.
\end{lemma}

\begin{proof}
Let $A_t=\|w_t-u\|_2^2$. Projection, convexity, and $\|g_t\|_2\le G$ imply
\begin{equation}
f_t(w_t)-f_t(u)\le \frac{A_t-A_{t+1}}{2\eta_t}+\frac{\eta_tG^2}{2}.
\label{eq:onestep}
\end{equation}
Summing from $s$ to $T$ gives
\begin{align}
\sum_{t=s}^T(f_t(w_t)-f_t(u))
&\le \frac{A_s}{2\eta_s}+\frac12\sum_{t=s+1}^TA_t\left(\frac1{\eta_t}-\frac1{\eta_{t-1}}\right)-\frac{A_{T+1}}{2\eta_T}+\frac{G^2}{2}\sum_{t=s}^T\eta_t.\label{eq:telescope}
\end{align}
The schedule $\eta_t$ is nondecreasing, so the middle coefficients and the final distance term are nonpositive. Write $\ell=T-s+1$. Since $A_s\le D^2$ and $\sum_{j=1}^{\ell}j^{-1/2}\le2\sqrt\ell$,
\[
\sum_{t=s}^T(f_t(w_t)-f_t(u))\le\frac{D^2}{2\eta_s}+\frac{G^2}{2}\sum_{t=s}^T\eta_t\le\frac{DG\sqrt\ell}{\sqrt2}+\frac{DG\sqrt\ell}{\sqrt2}=\sqrt2DG\sqrt\ell.
\]
\end{proof}

More generally, the schedule $\eta_t=cD/(G\sqrt{T-t+1})$ gives suffix constant $c+1/(2c)$. The choice $c=1/\sqrt2$ in Algorithm~\ref{alg:deadline} minimizes this expression and yields $\sqrt2$.

The discounted guarantee now follows from an exact identity.

\begin{lemma}[Geometric suffix mixture]
\label{lem:abel}
For any real sequence $r_1,\ldots,r_T$, define $S_j=\sum_{t=T-j+1}^Tr_t$. For every $\lambda\in[0,1]$,
\begin{equation}
\sum_{t=1}^T\lambda^{T-t}r_t=(1-\lambda)\sum_{j=1}^{T-1}\lambda^{j-1}S_j+\lambda^{T-1}S_T.
\label{eq:abel}
\end{equation}
The coefficients on the right are nonnegative, sum to one, and have mean suffix length $H_T(\lambda)$.
\end{lemma}

\begin{proof}
The coefficient of $r_{T-k+1}$ on the right is $(1-\lambda)\sum_{j=k}^{T-1}\lambda^{j-1}+\lambda^{T-1}=\lambda^{k-1}$. This proves the identity. The coefficient sum is one. If $J$ has these coefficients as its distribution on $[T]$, the tail-sum formula gives $\E[J]=\sum_{j=1}^T\Pr(J\ge j)=\sum_{j=1}^T\lambda^{j-1}=H_T(\lambda)$.
\end{proof}

\begin{theorem}[Uniform discount adaptation]
\label{thm:upper}
Under the assumptions of Section~\ref{sec:setting}, Algorithm~\ref{alg:deadline} satisfies, on every loss sequence,
\begin{equation}
\forall\lambda\in[0,1],\ \forall u\in\W,\qquad
R_T(\lambda,u)\le\sqrt2DG\sqrt{H_T(\lambda)}.
\label{eq:upper}
\end{equation}
Consequently, the same bound holds for $\DReg_T(\lambda)$ simultaneously for every $\lambda\in[0,1]$.
\end{theorem}

\begin{proof}
Fix $u$ and apply Lemma~\ref{lem:abel} to $r_t=f_t(w_t)-f_t(u)$. Lemma~\ref{lem:suffix} bounds $S_j\le\sqrt2DG\sqrt j$ for every $j$. With the random suffix length $J$ from Lemma~\ref{lem:abel},
\[
R_T(\lambda,u)\le\sqrt2DG\E[\sqrt J]\le\sqrt2DG\sqrt{\E[J]}=\sqrt2DG\sqrt{H_T(\lambda)},
\]
where the second inequality is Jensen's inequality. No step depends on $\lambda$, so the statement holds for the continuum simultaneously. Maximizing over $u$ proves the discounted-regret claim.
\end{proof}

The finite-horizon form matters near $\lambda=1$. Since
\[
H_T(\lambda)\le\min\left\{T,\frac1{1-\lambda}\right\},
\]
Theorem~\ref{thm:upper} implies the familiar $O(DG/\sqrt{1-\lambda})$ rate and never exceeds ordinary $O(DG\sqrt T)$ regret. At $\lambda=0$, only the last loss is evaluated and the bound is $\sqrt2DG$.

\section{The rate is pointwise minimax}
\label{sec:lower}

The next theorem grants the learner the discount. It therefore lower bounds both factor-aware algorithms and any algorithm required to work uniformly.

\begin{theorem}[Oracle lower bound]
\label{thm:lower}
Fix $T\ge1$ and $\lambda\in[0,1]$. On $\W=[-D/2,D/2]$, for every possibly randomized online algorithm, there is a distribution over convex losses with subgradients bounded by $G$ such that
\begin{equation}
\E[\DReg_T(\lambda)]\ge\frac14DG\sqrt{H_T(\lambda)}.
\label{eq:lower}
\end{equation}
The result continues to hold if the algorithm knows $\lambda$ before play.
\end{theorem}

\begin{proof}
Let $\varepsilon_1,\ldots,\varepsilon_T$ be independent Rademacher variables and set $f_t(w)=G\varepsilon_tw$. Adding the constant $GD/2$ to every $f_t$ places all losses in $[0,GD]$ without changing regret. A learner's $w_t$ is measurable with respect to its internal randomness and $\varepsilon_1,\ldots,\varepsilon_{t-1}$. Thus its expected weighted loss is zero before the harmless shift. The best comparator has weighted loss
\[
\min_{u\in[-D/2,D/2]}Gu\sum_{t=1}^T\lambda^{T-t}\varepsilon_t=-\frac{DG}{2}\left|\sum_{t=1}^T\lambda^{T-t}\varepsilon_t\right|.
\]
The sharp $p=1$ Khintchine inequality~\citep{haagerup1981best} and independence yield
\begin{align}
\E[\DReg_T(\lambda)]
&\ge\frac{DG}{2\sqrt2}\sqrt{\sum_{j=0}^{T-1}\lambda^{2j}}.
\end{align}
Finally,
\[
\sum_{j=0}^{T-1}\lambda^{2j}=H_T(\lambda)\frac{1+\lambda^T}{1+\lambda}\ge\frac12H_T(\lambda).
\]
Combining the two displays proves~\eqref{eq:lower}. The argument conditions on any internal random seed, so it applies to randomized algorithms and to algorithms that use $\lambda$.
\end{proof}

\begin{corollary}[Zero adaptivity price]
\label{cor:minimax}
Let $V_T(\lambda)$ be the minimax expected discounted regret when $\lambda$ is revealed before play. For every $\lambda\in[0,1]$,
\[
\frac14DG\sqrt{H_T(\lambda)}\le V_T(\lambda)\le\sqrt2DG\sqrt{H_T(\lambda)}.
\]
Algorithm~\ref{alg:deadline}, which is not told $\lambda$, attains the upper bound for all $\lambda$ simultaneously.
\end{corollary}

This is stronger than selecting one tuning after the fact. A single sequence $(w_t)_{t=1}^T$ competes, within a universal constant of the oracle minimax value, under uncountably many objectives. No union bound, discretization, or confidence parameter appears because Theorem~\ref{thm:upper} is deterministic.

\section{Terminal kernels and mirror geometry}
\label{sec:general}

The argument exposes a more general reduction. Index a terminal weighting kernel by age, so $a_1$ multiplies the newest loss and $a_T$ the oldest. Assume
\begin{equation}
1=a_1\ge a_2\ge\cdots\ge a_T\ge0,
\label{eq:kernel}
\end{equation}
and set $a_{T+1}=0$. Its effective mass is $H(a)=\sum_{j=1}^Ta_j$.

\begin{theorem}[Terminal-kernel guarantee]
\label{thm:kernel}
Algorithm~\ref{alg:deadline} satisfies simultaneously for every kernel obeying~\eqref{eq:kernel} and every $u\in\W$,
\begin{equation}
\sum_{j=1}^Ta_j\bigl(f_{T-j+1}(w_{T-j+1})-f_{T-j+1}(u)\bigr)
\le\sqrt2DG\sqrt{H(a)}.
\label{eq:kernelbound}
\end{equation}
\end{theorem}

\begin{proof}
Let $d_j=a_j-a_{j+1}$. Then $d_j\ge0$, $\sum_jd_j=a_1=1$, and discrete Abel summation gives the weighted regret as $\sum_{j=1}^Td_jS_j$. Lemma~\ref{lem:suffix} and Jensen give at most $\sqrt2DG\sqrt{\sum_jjd_j}$. A second summation by parts gives $\sum_jjd_j=\sum_ja_j=H(a)$.
\end{proof}

Thus the kernel may also be chosen after the losses. Exponential discounting takes $a_j=\lambda^{j-1}$ and recovers Theorem~\ref{thm:upper}. The next result identifies the oracle scale for an arbitrary fixed kernel.

\begin{proposition}[Fixed-kernel lower bound]
\label{prop:kernellower}
Fix any nonnegative terminal weights $a_1,\ldots,a_T$. Even if these weights are revealed before play, every online algorithm has worst-case expected weighted regret at least
\[
\frac{DG}{2\sqrt2}\left(\sum_{j=1}^Ta_j^2\right)^{1/2}.
\]
\end{proposition}

\begin{proof}
Use the Rademacher linear losses from Theorem~\ref{thm:lower}, now with coefficients $a_j$. The learner's expected loss is zero, the best comparator contributes $(DG/2)|\sum_ja_j\varepsilon_j|$, and Khintchine's inequality completes the proof.
\end{proof}

The upper bound in Theorem~\ref{thm:kernel} depends on $\|a\|_1^{1/2}$, while Proposition~\ref{prop:kernellower} depends on $\|a\|_2$. These scales are equivalent up to a constant for geometric kernels because $H_T(\lambda)\le2\sum_{j=0}^{T-1}\lambda^{2j}$. This is why Corollary~\ref{cor:minimax} is sharp. They need not be equivalent for every slowly decaying kernel, so we do not claim pointwise minimaxity for all of~\eqref{eq:kernel}.

The Euclidean update is also inessential. Let $\|\cdot\|$ be a norm with dual norm $\|\cdot\|_*$, and let $\psi$ be one-strongly convex on $\W$. Assume its Bregman divergence satisfies
\[
\sup_{u,w\in\W}D_\psi(u,w)\le\frac{B^2}{2}.
\]
Online mirror descent uses
\begin{equation}
w_{t+1}\in\arg\min_{w\in\W}\left\{\eta_t\langle g_t,w\rangle+D_\psi(w,w_t)\right\},
\qquad
\eta_t=\frac{B}{G\sqrt{2(T-t+1)}},
\label{eq:omd}
\end{equation}
where $\|g_t\|_*\le G$.

\begin{corollary}[Deadline mirror descent]
\label{cor:mirror}
The suffix, discount, and terminal-kernel results hold for~\eqref{eq:omd} with $D$ replaced by $B$.
\end{corollary}

The proof is in Appendix~\ref{app:mirror}. The standard mirror-descent inequality has the same potential difference as~\eqref{eq:onestep}. Its interior telescoping terms again have the favorable sign because $\eta_t$ increases.

\section{Relation to prior work and scope}
\label{sec:related}

\paragraph{Discounted online learning.} Freund and Hsu~\citep{freund2008hedging} and Kapralov and Panigrahy~\citep{kapralov2011prediction} developed discounted predictors for expert advice. Kapralov and Panigrahy also outlined guarantees uniform over time scales by combining discounted predictors. Those guarantees carry logarithmic scale or confidence terms and address expert aggregation. Chernov and Zhdanov~\citep{chernov2010prediction} studied discounted expert loss through defensive forecasting. More recent work connects discounting to dynamic regret and adaptive optimization~\citep{jacobsen2024online,ahn2024understanding,xie2026dynamic}.

Zhang et al.~\citep{zhang2024discounted} study instance-adaptive FTRL for supplied, possibly time-varying factors. For a constant known factor, they recover the minimax scale through an effective gradient variance. They explicitly leave online selection of the factor open. Yang et al.~\citep{yang2026uniform} address an unknown fixed factor and obtain the first uniform OCO result over a continuum. Their algorithm needs the terminal $T$ to construct its factor grid, so the known-deadline assumption does not explain the gap between their rate and ours. Our gain instead comes from targeting terminal suffixes directly rather than combining factor-specific learners.

\paragraph{Interval adaptivity.} Adaptive and strongly adaptive regret require performance on intervals with arbitrary left and right endpoints~\citep{daniely2015strongly,cutkosky2020parameter}. Their algorithms typically maintain learners started at different times and pay logarithmic overhead. Deadline OGD controls the nested family $[s,T]$ for one fixed right endpoint. This smaller family is sufficient because every nonincreasing terminal kernel is a mixture of its indicators. Requiring the same schedule to work at every stopping time would restore all interval endpoints and is outside our theorem.

\paragraph{Knowledge assumptions.} The algorithm uses upper bounds $D$ and $G$ and the deadline $T$. It uses neither the discount nor the loss range. Unknown scale can be addressed by standard parameter-free machinery, but combining that machinery with the exact terminal profile is a separate question. The guarantee is adversarial and deterministic. It therefore also applies to stochastic losses, discounts chosen after the data, and adaptive environments that respect the online protocol.

\paragraph{Why the simple schedule was missed.} Known-factor analyses use a constant learning rate matched to the effective horizon. Unknown-parameter methods then naturally aggregate these constant-rate experts. The alternative is to match the remaining time, not the unknown effective time. Increasing steps look unfavorable for ordinary prefix regret, but their potential correction has the correct sign on every terminal suffix. Abel summation transfers precisely that one-sided guarantee to every discount.

\section*{Reproducibility statement}

The assumptions and full algorithm are stated in Sections~\ref{sec:setting} and~\ref{sec:upper}. The main upper and lower bounds include complete proofs in Sections~\ref{sec:upper} and~\ref{sec:lower}. The appendix gives the mirror-descent proof, the general terminal-mixture reduction, lower-bound constants, and endpoint checks. This theoretical work uses no datasets, experiments, or learned artifacts.

\bibliography{references}
\bibliographystyle{iclr2027_conference}

\appendix

\section{Mirror-descent proof}
\label{app:mirror}

We prove Corollary~\ref{cor:mirror}. The standard online mirror-descent inequality states that for every $u\in\W$,
\begin{equation}
f_t(w_t)-f_t(u)\le\langle g_t,w_t-u\rangle
\le\frac{D_\psi(u,w_t)-D_\psi(u,w_{t+1})}{\eta_t}+\frac{\eta_t}{2}\|g_t\|_*^2.
\label{eq:mirrorone}
\end{equation}
For completeness, the optimality condition for~\eqref{eq:omd} and the three-point Bregman identity imply
\[
\eta_t\langle g_t,w_{t+1}-u\rangle\le D_\psi(u,w_t)-D_\psi(u,w_{t+1})-D_\psi(w_{t+1},w_t).
\]
Add $\eta_t\langle g_t,w_t-w_{t+1}\rangle$. One-strong convexity gives $D_\psi(w_{t+1},w_t)\ge\|w_{t+1}-w_t\|^2/2$, and Fenchel--Young gives
\[
\eta_t\langle g_t,w_t-w_{t+1}\rangle-\frac12\|w_t-w_{t+1}\|^2\le\frac{\eta_t^2}{2}\|g_t\|_*^2.
\]
Dividing by $\eta_t$ proves~\eqref{eq:mirrorone}.

Write $V_t=D_\psi(u,w_t)$. Summing~\eqref{eq:mirrorone} from $s$ through $T$ gives
\begin{align}
\sum_{t=s}^T(f_t(w_t)-f_t(u))
&\le\frac{V_s}{\eta_s}+\sum_{t=s+1}^TV_t\left(\frac1{\eta_t}-\frac1{\eta_{t-1}}\right)-\frac{V_{T+1}}{\eta_T}+\frac{G^2}{2}\sum_{t=s}^T\eta_t.
\end{align}
Both terms after $V_s/\eta_s$ involving divergences are nonpositive because $V_t\ge0$ and $\eta_t$ is nondecreasing. With $\ell=T-s+1$, the Bregman-diameter assumption and the same inverse-square-root sum used in Lemma~\ref{lem:suffix} give
\[
\sum_{t=s}^T(f_t(w_t)-f_t(u))\le\frac{B^2}{2\eta_s}+\frac{G^2}{2}\sum_{t=s}^T\eta_t\le\sqrt2BG\sqrt\ell.
\]
Lemmas~\ref{lem:abel} and Theorem~\ref{thm:kernel} are algebraic and therefore apply unchanged.

\section{A general terminal-mixture principle}
\label{app:mixture}

The proof can be separated completely from OGD. This abstraction may be useful when a problem admits a sharper terminal profile than $\sqrt j$.

\begin{proposition}[Suffix profile to kernel profile]
\label{prop:profile}
Fix a decision sequence and comparator $u$. Suppose its terminal suffix sums satisfy $S_j(u)\le B(j,u)$ for every $j\in[T]$. For any kernel $a$ satisfying~\eqref{eq:kernel},
\begin{equation}
\sum_{j=1}^Ta_jr_{T-j+1}(u)\le\sum_{j=1}^T(a_j-a_{j+1})B(j,u).
\label{eq:profile}
\end{equation}
If $B(\cdot,u)$ is nondecreasing and concave, the right side is at most $B(H(a),u)$ after extending $B$ to $[1,T]$ by linear interpolation.
\end{proposition}

\begin{proof}
The first claim is discrete Abel summation. The differences $d_j=a_j-a_{j+1}$ form a probability distribution. The second claim follows from Jensen and $\sum_jjd_j=H(a)$.
\end{proof}

The normalization $a_1=1$ is only for convenience. If $a_1=c>0$, divide the kernel by $c$ and multiply the resulting bound by $c$. If $a_1=0$, monotonicity makes all weights zero.

For the exponential kernel, Proposition~\ref{prop:profile} gives a probabilistic interpretation of discounting. Draw a random length $J$ from a geometric distribution with success probability $1-\lambda$, truncated by moving all tail mass to $T$. The discounted regret is exactly the expected regret on the last $J$ rounds. This identity concerns signed regrets, not absolute losses, and therefore preserves cancellation within each suffix.

\section{Lower-bound details and constants}
\label{app:lower}

We record three details about Theorem~\ref{thm:lower}. First, the loss distribution is oblivious. It therefore lower bounds algorithms against both oblivious and adaptive loss classes. Second, the learner may randomize. Conditional on its random seed and the past, $\varepsilon_t$ has zero mean, so
\[
\E\left[\sum_{t=1}^T\lambda^{T-t}G\varepsilon_tw_t\right]=0.
\]
Third, the expectation lower bound implies a worst-case sequence lower bound for expected regret by averaging over the Rademacher distribution.

The only external inequality used is
\[
\E\left|\sum_{j=1}^Tc_j\varepsilon_j\right|\ge\frac1{\sqrt2}\left(\sum_{j=1}^Tc_j^2\right)^{1/2}.
\]
With $c_j=\lambda^{j-1}$, its squared norm is
\[
Q_T(\lambda)=\sum_{j=0}^{T-1}\lambda^{2j}=
\begin{cases}
(1-\lambda^{2T})/(1-\lambda^2),&\lambda<1,\\
T,&\lambda=1.
\end{cases}
\]
For every $\lambda\in[0,1]$,
\[
\frac{Q_T(\lambda)}{H_T(\lambda)}=\frac{1+\lambda^T}{1+\lambda}\in\left[\frac12,2\right].
\]
In fact the ratio is at most one for $T\ge1$, but the lower inequality is the direction needed in the theorem. The constants then combine as
\[
\frac{DG}{2}\cdot\frac1{\sqrt2}\cdot\sqrt{\frac{H_T(\lambda)}2}=\frac14DG\sqrt{H_T(\lambda)}.
\]

\section{Endpoint and stopping-time checks}
\label{app:endpoints}

At $\lambda=0$, the identity~\eqref{eq:abel} reduces to $r_T=S_1$. Algorithm~\ref{alg:deadline} uses $\eta_T=D/(G\sqrt2)$ for the final update, although that update cannot affect the final loss. The suffix bound at length one comes from the potential already accumulated before the last prediction and remains valid.

At $\lambda=1$,~\eqref{eq:abel} reduces to $\sum_{t=1}^Tr_t=S_T$. The first step size is $D/(G\sqrt{2T})$, which optimally balances the two terms in our inverse-square-root-sum bound. Intermediate steps grow, but their potential corrections in~\eqref{eq:telescope} are nonpositive.

The schedule is not anytime. If play unexpectedly stops at $\tau<T$, then its early step sizes were calibrated to the later deadline $T$. One can restart on doubling epochs to obtain an anytime method, but the terminal weights would cross epoch boundaries and the exact constant profile need not survive. Theorem~\ref{thm:upper} deliberately addresses the same known-terminal setup in which the recent uniform-discount construction sets a grid using $T$~\citep{yang2026uniform}.

\end{document}
